\section{Resultados e Discuss\~ao}
\label{sec:resultados}

\textbf{UDVs.} Das 2.203 opiniões, 2.105 recebem evidência localizada na fala da própria pessoa, 277
por citação literal e 1.828 por maior cosseno (Tabela~\ref{tab:udv-niveis}), e o verificador não
encontrou inconsistência em nenhum registro. A camada semântica é a menos confiável: nas 183
opiniões de treino com citação confiável, a sentença de maior cosseno com a opinião completa é a da
própria citação em 112 (61,2\%, IC 95\% de 53,0\% a 68,8\%), e, com os trechos entre aspas removidos,
em 12 de 142 consultas do treino e em 4 de 40 da validação. Esses números indicam que a sentença
escolhida por similaridade difere com frequência da passagem citada, mas não separam o erro da
escolha de outra sentença que também sustente a opinião; essa separação cabe à validação humana, na
qual <resultado da validação humana: precisão estrita da citação literal e tolerante dos casos
semânticos de alta confiança, com intervalos de Wilson, critérios atingidos ou não, e concordância do
anotador consigo mesmo>.

\textbf{Simulação.} Na validação, <n> das <n> opiniões ligadas a atores com perfil geraram perguntas,
e o maior acerto foi obtido com $k = {<}k{>}$ na condição 2 e com $\gamma = {<}\gamma{>}$ na condição 3.
No teste, <n> das 106 opiniões ligadas geraram perguntas, de <n> atores em <n> audiências, e <n>
foram descartadas por falta de três outros participantes (Tabela~\ref{tab:resultados}). A soma das
probabilidades das quatro letras teve média <0,00> e mínimo <0,00> nas condições 0 a 2. Na condição
2, as perguntas classificadas como DIRETA, INDIRETA e ESPECULATIVA foram <n>, <n> e <n>, com acerto
de <0,00>, <0,00> e <0,00>, e <n> respostas não trouxeram rótulo. Nos <n> pedidos da geração aberta,
as condições 1, 2 e 3 produziram <n>, <n> e <n> recusas, e a fração de frases sem fundamento na
conferência da justificação foi <0,00>, <0,00> e <0,00>. <Discussão dos resultados da simulação, a
escrever com os números da execução completa.>

\input{tables/main-results}

\textbf{Ameaças à validade.} A opção correta de cada pergunta é a opinião que a matéria atribui à
pessoa, na redação da matéria, e o acerto mede a capacidade de reconhecer essa atribuição, que é uma
seleção editorial do que foi dito. Os distratores são opiniões de outros participantes da mesma
audiência e podem coincidir com posições que a pessoa também sustenta; essas perguntas admitem mais
de uma resposta plausível e reduzem o acerto máximo possível, e sua frequência não foi medida. O nome
da pessoa está presente em todas as condições, e o modelo pode associar a ele posições conhecidas
fora do conjunto; a condição 0 mede esse efeito, que as diferenças entre condições só cancelam se
for aditivo. Quando o cargo registrado traz a função da pessoa na própria audiência de teste, como
relator ou presidente da comissão, essa informação entra em todas as condições, eleva o acerto
absoluto e se cancela nas diferenças. As perguntas de teste são no máximo 106, concentradas em 27
audiências, e só as diferenças cujo intervalo exclui o zero são tratadas como efeito. Os valores de
$k$ e de $\gamma$ são escolhidos para o conjunto das perguntas de validação, e não por ator. A
fidelidade dos perfis às falas não foi medida por anotação humana; as regras do prompt resultam da
leitura de perfis gerados durante o desenvolvimento (Seção~\ref{sec:metodo-perfis}). Por fim, a
conferência da justificação comprova a existência dos itens citados e a literalidade das cópias, mas
não a sustentação das frases.

\textbf{Uso das simulações.} A fala simulada é uma estimativa condicionada às falas públicas da
pessoa nas audiências de treino, e não uma declaração dela. Por isso, cada saída acompanha a base da
estimativa e a justificação conferida, o partido não entra no prompt, e a recusa é uma saída prevista
quando o material não se aplica. As simulações não devem ser apresentadas como citação nem usadas
para inferir intenção política.
